\documentclass[addpoints,spanish, 12pt,a4paper]{exam}
%\documentclass[answers, spanish, 12pt,a4paper]{exam}
% \printanswers
\renewcommand*\half{.5}
\pointpoints{punto}{puntos}
\hpword{Puntos:}
\vpword{Puntos:}
\htword{Total}
\vtword{Total}
\hsword{Resultado:}
\hqword{Ejercicio:}
\vqword{Ejercicio:}

\usepackage[utf8]{inputenc}
\usepackage[spanish]{babel}
\usepackage{eurosym}
%\usepackage[spanish,es-lcroman, es-tabla, es-noshorthands]{babel}

\usepackage[margin=1in]{geometry}
\usepackage{amsmath,amssymb}
\usepackage{multicol}
\usepackage{yhmath}

\pointsinrightmargin % Para poner las puntuaciones a la derecha. Se puede cambiar. Si se comenta, sale a la izquierda.
\extrawidth{-2.4cm} %Un poquito más de margen por si ponemos textos largos.
\marginpointname{ \emph{\points}}

\usepackage{graphicx}

\graphicspath{{../../img/}}

\newcommand{\class}{2ª Bachillerato Sociales}
\newcommand{\examdate}{\today}
\newcommand{\examnum}{Probabilidad}
\newcommand{\tipo}{A}

\newcommand{\timelimit}{45 minutos}

\renewcommand{\solutiontitle}{\noindent\textbf{Solución:}\enspace}

\pagestyle{head}
\firstpageheader{\includegraphics[width=0.2\columnwidth]{header_left}}{\textbf{Departamento de Matemáticas\linebreak \class}\linebreak \examnum}{\includegraphics[width=0.1\columnwidth]{header_right}}
\runningheader{\class}{\examnum}{Página \thepage\ de \numpages}
\runningheadrule

\usepackage{pgf,tikz,pgfplots}
\pgfplotsset{compat=1.15}
\usepackage{mathrsfs}
\usetikzlibrary{arrows}

\begin{document}

\noindent
\begin{tabular*}{\textwidth}{l @{\extracolsep{\fill}} r @{\extracolsep{6pt}} }
\textbf{Nombre:} \makebox[3.5in]{\hrulefill} & \textbf{Fecha:}\makebox[1in]{\hrulefill} \\
 & \\
\textbf{Tiempo: \timelimit} & Tipo: \tipo 
\end{tabular*}
\rule[2ex]{\textwidth}{2pt}
Esta prueba tiene \numquestions\ ejercicios. La puntuación máxima es de \numpoints.
La nota final de la prueba será la parte proporcional de la puntuación obtenida sobre la puntuación máxima.

\begin{center}

\addpoints
 %\gradetable[h][questions]
\pointtable[h][questions]
\end{center}

\noindent
\rule[2ex]{\textwidth}{2pt}

\begin{questions}

\question[3] De los 39 alumnos de una clase,
16 escogieron francés, como idioma y 27 inglés.
9 eligieron ambos idiomas y el resto no optó por ninguno de ellos. elegido un alumno al azar, calcula las probabilidades de que escogiera:
\begin{multicols}{2}
\begin{parts}
    \part Francés
    \part Inglés
    \part Ambos
    \part Francés o Inglés
    \part Francés, pero no inglés
    \part Inglés, pero no francés
\end{parts}
\end{multicols}
\begin{solution}
    16/39, 9/13 , 3/13 , 34/39, 7/39, 6/13
\end{solution}


\question En un grupo de estudiantes, un 10\% sabe inglés y alemán, un 50\% sabe inglés pero no alemán y, entre los que saben alemán, un 40\% sabe inglés.
\begin{parts}
  \part[1] ¿Qué porcentaje de estudiantes sabe inglés?
    \begin{solution}
     $ \begin{matrix} 
           & I & \overline{I} &  \\
          A &10 & & \\
          \overline{A} &50 & & \\
                       &60&40&100
      \end{matrix} $
      \\
      $50+10=60$\% de los estudiantes sabe inglés.

    \end{solution}
  \part[2] ¿Qué porcentaje sabe alemán?
  \begin{solution}$P(I|A)=0.4 \to \dfrac{P(I\cap A)}{P(A)}=0.4\to P(A)=\dfrac{0.1}{0.4}=\dfrac{1}{4}=0.25=25$\% sabe alemán \\
     $ \begin{matrix} 
           & I & \overline{I} &  \\
          A &10 & 15& 25\\
          \overline{A} &50 &25 &75 \\
                       &60&40&100
      \end{matrix} $
  \end{solution}
  

  \part[1] ¿Qué porcentaje sabe alguno de los dos idiomas?
  \begin{solution}
    $1-P(Ninguno)=1-0.25=0.75=75$\% sabe alguno
  \end{solution}
\end{parts}

\question[2] Dados dos sucesos A y B de un mismo experimento aleatorio, se sabe que 
  $P(A)=0.5$ , $P(B)=0.3$ y $P(A\cup B)=0.6$. 
  ¿Son A y B independientes?
  \begin{solution}
  $P(A\cup B)=P(A)+P(B)-P(A\cap B) 
  \to P(A\cap B)=P(A)+P(B)-P(A\cup B)=0.5+0.3-0.6=0.2\neq 0.15=0.5 \cdot 0.3=P(A) \cdot P(B)$ \\
  Por tanto no son independientes  
  \end{solution}

%\question La probabilidad de que un autobús llegue con retraso a una parada es $0.2$.  Si pasa cuatro veces a lo largo del día por la parada, calcular la probabilidad de que:  
% \begin{parts}
%   \part[0\half]  No llegue con retraso ninguna de las veces.
%     \begin{solution}
%         \[
%         P(X = 0) = (0.8)^4 = \mathbf{0.4096}
%         \]
%     \end{solution}

%     \part[0\half] Llegue con retraso al menos una vez.
%     \begin{solution}
%         \[
%         P(X \geq 1) = 1 - P(X = 0) = \mathbf{0.5904}
%         \]
%     \end{solution}

%     \part[0\half] Al menos tres veces llegue con retraso.
%     \begin{solution}
%         \[
%         P(X \geq 3) = P(X = 3) + P(X = 4) = \mathbf{0.0272}
%         \]
%     \end{solution}

%     \part[0\half] Llegue con retraso exactamente dos veces consecutivas.
%     \begin{solution}
%         \[
%         P(\text{exactamente 2 retrasos consecutivos}) = \mathbf{0.0768}
%         \]
%     \end{solution}
% \end{parts}

% \question Para efectuar cierto diagnóstico, un hospital dispone de dos escáneres, A y B.
% el 65\% de las pruebas se realizan con el escáner A y el resto con B.
% Se sabe que el escáner A tiene un 5\% de diagnósticos erróneos, mientras que el escáner B tiene un 8\% de diagnósticos erróneos.

% Calcular la probabilidad de que:

% \begin{parts}
%   \part[1] El diagnóstico efectuado a un paciente en ese hospital sea erróneo.
%     \begin{solution}
%         \[
%         P(E) = (0.05 \times 0.65) + (0.08 \times 0.35) = 0.0325 + 0.028 = \mathbf{0.0605}
%         \]
%     \end{solution}

%     \part[2] El diagnóstico se haya efectuado usando el escáner A, sabiendo que ha resultado erróneo.
%     \begin{solution}
%         \[
%         P(A | E) = \frac{0.05 \times 0.65}{0.0605} = \frac{0.0325}{0.0605} \approx \mathbf{0.5376}
%         \]
%     \end{solution}
% \end{parts}

% \question Los escolares de un cierto colegio de Zaragoza fueron encuestados acerca de su alimentación y ejercicio físico.  
% Una proporción de $\frac{2}{5}$ hacían ejercicio regularmente y $\frac{2}{3}$ siempre desayunaban.  
% Además, entre los que siempre desayunan, una proporción de $\frac{9}{25}$ hacían ejercicio regularmente.  
% Se elige al azar un escolar de ese colegio.  

% \begin{parts}
%     \part[1] ¿Es independiente que siempre desayune y que haga ejercicio regularmente?
%     \begin{solution}
%         \[
%         \begin{array}{c|cc|c}
%              & D & \neg D & \text{Total} \\
%             \hline
%             E & \frac{6}{25} & 4/25 & \frac{2}{5} \\
%             \neg E & 32/75 & 13/75 & 3/5 \\
%             \hline
%             \text{Total} & \frac{2}{3} & 1/3 & 1
%         \end{array}
%         \]
%         Definimos los eventos:  
%         - \( D \): Siempre desayuna.  
%         - \( E \): Hace ejercicio regularmente.  

%         Probabilidades dadas:  
%         \[
%         P(E) = \frac{2}{5}, \quad P(D) = \frac{2}{3}, \quad P(E | D) = \frac{9}{25}
%         \]
%         Calculamos \( P(E) \cdot P(D) \):
%         \[
%         P(E) \cdot P(D) = \left(\frac{2}{5} \right) \left(\frac{2}{3} \right) = \frac{4}{15}
%         \]
%         Calculamos \( P(E \cap D) \):
%         \[
%         P(E \cap D) = P(E | D) P(D) = \left(\frac{9}{25} \right) \left(\frac{2}{3} \right) = \frac{18}{75} = \frac{6}{25}
%         \]

%         Como \( P(E \cap D) \neq P(E) \cdot P(D) \), los eventos **no son independientes**.  
%         \[
%         \mathbf{No\ son\ independientes.}
%         \]
%     \end{solution}

%     \part[2] Calcular la probabilidad de que no siempre desayune y no haga ejercicio regularmente.
%     \begin{solution}
%         Probabilidad de que no siempre desayune ni haga ejercicio:
%         \[
%         P(\neg D \cap \neg E) = \dfrac{13}{75}
%         \]

%     \end{solution}
% \end{parts}



\question Una empresa fabrica dos modelos de ordenadores portátiles, siendo la producción del modelo
B la tercera parte de todos los que se fabrican. Se sabe que la probabilidad de que un ordenador portátil del
modelo A salga defectuoso es de 0,02, mientras que esa probabilidad en el modelo B es de 0,06.

\begin{parts}
    \part[1] Calcula la probabilidad de que un ordenador fabricado por dicha empresa elegido al azar no salga defectuoso.
    \begin{solution}
        Sabemos que:
        - La probabilidad de que un ordenador del **modelo A** salga defectuoso es \( P(D|A) = 0.02 \),
        - La probabilidad de que un ordenador del **modelo B** salga defectuoso es \( P(D|B) = 0.06 \).

        Además, se nos dice que la producción del **modelo B** es la tercera parte de todos los ordenadores fabricados, es decir:
        - La probabilidad de que un ordenador sea del **modelo A** es \( P(A) = \frac{2}{3} \),
        - La probabilidad de que un ordenador sea del **modelo B** es \( P(B) = \frac{1}{3} \).

        La probabilidad de que un ordenador **no sea defectuoso** es:
        \[
        P(\text{no defectuoso}) = P(\text{no defectuoso} | A) \cdot P(A) + P(\text{no defectuoso} | B) \cdot P(B)
        \]
        Sabemos que:
        - \( P(\text{no defectuoso} | A) = 1 - P(D|A) = 0.98 \),
        - \( P(\text{no defectuoso} | B) = 1 - P(D|B) = 0.94 \).

        Sustituyendo:
        \[
        P(\text{no defectuoso}) = 0.98 \cdot \frac{2}{3} + 0.94 \cdot \frac{1}{3} 
        \]
        \[
        P(\text{no defectuoso}) = 0.6533 + 0.3133 = 0.9666
        \]

        Por lo tanto, la probabilidad de que un ordenador fabricado por la empresa elegido al azar no salga defectuoso es aproximadamente **0.967**.
    \end{solution}

    \part[2] Sea del modelo A, si se sabe que ha salido defectuoso.
    \begin{solution}
        Utilizamos el teorema de Bayes para calcular la probabilidad de que un ordenador sea del modelo A, dado que ha salido defectuoso:

        \[
        P(A | D) = \frac{P(D | A) \cdot P(A)}{P(D)}
        \]

        Sabemos que:
        - \( P(D | A) = 0.02 \),
        - \( P(A) = \frac{2}{3} \),
        - \( P(D | B) = 0.06 \),
        - \( P(B) = \frac{1}{3} \).

        Ya hemos calculado \( P(D) \), la probabilidad total de que un ordenador salga defectuoso:
        \[
        P(D) = P(D | A) \cdot P(A) + P(D | B) \cdot P(B)
        \]
        \[
        P(D) = 0.02 \cdot \frac{2}{3} + 0.06 \cdot \frac{1}{3}
        \]
        \[
        P(D) = 0.0133 + 0.02 = 0.0333
        \]

        Ahora, aplicamos el teorema de Bayes:
        \[
        P(A | D) = \frac{0.02 \cdot \frac{2}{3}}{0.0333} = \frac{0.0133}{0.0333} = 0.4
        \]

        Por lo tanto, la probabilidad de que el ordenador sea del modelo A, dado que ha salido defectuoso, es **0.4** o **40%**.
    \end{solution}

    \part[1] Aproximadamente, cuántos defectuosos habrá en un pedido de 1500 ordenadores.
    \begin{solution}
        Sabemos que la probabilidad de que un ordenador salga defectuoso es \( P(D) = 0.0333 \).

        El número esperado de defectuosos en un pedido de 1500 ordenadores es:
        \[
        \text{Defectuosos} = 1500 \times P(D) = 1500 \times 0.0333 = 50
        \]

        Por lo tanto, en un pedido de 1500 ordenadores, se espera que haya aproximadamente **50 defectuosos**.
    \end{solution}

\end{parts}



\end{questions}
\end{document}

